problem:  
Let \( (X^n, d) \) be a compact Alexandrov space of curvature \(\ge 0\) whose regular part \( M_{\mathrm{reg}} \subset X\) is a smooth Riemannian manifold of nonnegative sectional curvature, and assume the singular set \(S = X \setminus M_{\mathrm{reg}}\) has Hausdorff codimension at least \(3\). Suppose \(X\) carries an effective, isometric \(\mathrm{SO}(3)\)-action with 3‑dimensional principal orbits (so \(X/\!\!/\mathrm{SO}(3)\) is of dimension \(n-3\)).  

**Question:**  
Does such an Alexandrov space \(X\) necessarily admit a continuous family of smooth, closed Riemannian manifolds \( (M_i^n, g_i)\) with \(\sec(g_i)\ge 0\) converging to \(X\) in the Gromov–Hausdorff topology?  

Equivalently: does the presence of a 3‑dimensional isometric symmetry group and codimension‑\(\ge 3\) singularities force \(X\) to be *smoothable as a nonnegatively curved limit space*?

Why is it a "good" problem:  
This problem probes the fundamental interface between Alexandrov geometry, transformation group theory, and smoothing phenomena under curvature bounds.  

1. **Conceptual depth:**  
   The question asks when metric singularities in Alexandrov spaces cease to be genuine obstructions once curvature and symmetry constraints are present. Codimension‑\(\ge3\) singular sets are often topologically insignificant (as in manifold smoothing theory), but their geometric removability under nonnegative curvature remains mysterious. Determining whether such singularities can always be “smoothed away” in the Gromov–Hausdorff sense connects curvature comparison geometry with the topological smoothing theory of high-dimensional manifolds.

2. **Bridge between fields:**  
   The problem touches Alexandrov geometry (metric curvature bounds), differential geometry (nonnegative sectional curvature), transformation group actions (via isometric \(\mathrm{SO}(3)\)-symmetry and cohomogeneity structures), and metric limit theory (Gromov–Hausdorff convergence). A positive resolution would unify aspects of comparison geometry, topological transformation groups, and geometric analysis of limit spaces.

3. **Potential significance:**  
   If one could prove that symmetry and codimension‑3 singularities assure smoothability, it would establish a new *symmetric smoothing principle* for Alexandrov spaces, potentially extending Perelman’s stability theorem and the Alexandrov-to-Riemannian transition program. Conversely, a counterexample would exhibit new, symmetry-compatible nonsmooth curvature models, reshaping understanding of singularities in nonnegatively curved spaces.

4. **Simplicity and profundity:**  
   The statement is geometrically clear and concise: it asks whether symmetry and mild singularities suffice to recover manifold smoothability under a fundamental curvature constraint. Yet the resolution demands deep technical insight into the interplay of geometry, topology, and group actions—an archetype of “narrow entrance, profound depth.”